\section{条件异方差的统计特征}
\label{sec:11-Mathematical-Statistics/10-GARCH-and-Volatility/01}

你接手一串按天记录的收益率，第一件事是画自相关图。滞后一到二十，几乎全部落在显著性带里，看不出任何可预测的结构。你写下结论：这是白噪声，均值模型无从下手，交给下游按独立同分布处理。

三个月后，风控那边报过来一堆问题。同一个固定阈值的告警，在某几周每天响七八次，在另外几周整月一声不吭。你调高阈值，安静的时段更安静了，混乱的时段还是照样刷屏。这不像是阈值选得不对，倒像是被监控的东西在换性格。

问题出在你只检查了一阶。序列本身不相关，是关于 \(\Cov(r_t,r_{t-k})\) 的陈述；而告警刷屏关心的是 \(r_t\) 的幅度，那是二阶的事。这两件事可以完全脱钩：一个序列的一阶自相关全为零，它的平方序列却可以有强而持久的自相关。

\subsection{序列本身不相关，平方序列却相关}

把同一串数据做两遍自相关：一遍对 \(r_t\)，一遍对 \(r_t^2\)。前者是零，后者是一条缓慢下降、长期停在显著性带之外的曲线。这个反差是波动率聚集在统计上留下的指纹——大波动后面跟着大波动，平静后面跟着平静，幅度有记忆，方向没有。

\begin{center}
\begin{tikzpicture}[>=stealth]
% ---- 上：收益率序列 ----
\draw[gray!35] (0,0) -- (10.2,0);
\draw[->] (0,-1.75) -- (0,1.9) node[above,font=\scriptsize] {\(r_t\)};
\draw[->] (0,-1.75) -- (10.5,-1.75) node[right,font=\scriptsize] {\(t\)};
\draw[blue!60!black,line width=0.35pt] plot coordinates {
(0.000,-0.165) (0.105,0.049) (0.211,-0.472) (0.316,-0.087) (0.421,0.022) (0.526,0.022)
(0.632,-0.308) (0.737,0.051) (0.842,0.031) (0.947,-0.193) (1.053,0.148) (1.158,-0.010)
(1.263,0.003) (1.368,-0.080) (1.474,0.094) (1.579,0.016) (1.684,0.182) (1.789,0.075)
(1.895,-0.099) (2.000,-0.035) (2.105,0.140) (2.211,-0.057) (2.316,0.058) (2.421,0.077)
(2.526,-0.026) (2.632,-0.326) (2.737,0.023) (2.842,0.036) (2.947,0.022) (3.053,-0.117)
(3.158,0.187) (3.263,-0.022) (3.368,-0.097) (3.474,-0.110) (3.579,0.055) (3.684,-0.162)
(3.789,-0.153) (3.895,0.313) (4.000,0.241) (4.105,0.203) (4.211,0.267) (4.316,0.123)
(4.421,-0.337) (4.526,0.264) (4.632,0.281) (4.737,0.115) (4.842,-0.439) (4.947,0.321)
(5.053,0.362) (5.158,-0.137) (5.263,0.034) (5.368,0.193) (5.474,-0.021) (5.579,0.252)
(5.684,0.019) (5.789,0.152) (5.895,0.018) (6.000,-0.300) (6.105,-0.216) (6.211,0.690)
(6.316,0.087) (6.421,0.383) (6.526,0.344) (6.632,0.525) (6.737,0.187) (6.842,-0.191)
(6.947,0.006) (7.053,-0.552) (7.158,-0.062) (7.263,0.279) (7.368,-0.655) (7.474,0.071)
(7.579,0.277) (7.684,0.436) (7.789,-0.315) (7.895,-0.218) (8.000,-0.604) (8.105,-0.370)
(8.211,0.193) (8.316,0.712) (8.421,-0.460) (8.526,0.558) (8.632,0.180) (8.737,-0.197)
(8.842,0.827) (8.947,-1.132) (9.053,-0.038) (9.158,0.115) (9.263,1.047) (9.368,1.038)
(9.474,1.450) (9.579,0.103) (9.684,0.702) (9.789,0.963) (9.895,0.375) (10.000,0.249)};
\draw[gray!55,dashed] (6.05,-1.75) -- (6.05,1.75);
\node[font=\scriptsize,gray!45!black] at (2.9,1.62) {平静段};
\node[font=\scriptsize,gray!45!black] at (8.2,1.62) {高波动段};
\draw[gray!55,dashed] (0,0.9) -- (10.2,0.9);
\draw[gray!55,dashed] (0,-0.9) -- (10.2,-0.9);
\node[font=\scriptsize,gray!45!black,right] at (10.2,0.9) {\(\pm c\)};

% ---- 下左：r 的自相关 ----
\begin{scope}[shift={(0,-4.3)}]
\fill[gray!14] (0,-0.24) rectangle (4.3,0.24);
\draw[->] (0,0) -- (4.6,0) node[right,font=\scriptsize] {\(k\)};
\draw[->] (0,-0.45) -- (0,1.55);
\foreach \k/\h in {1/0.12,2/-0.16,3/0.04,4/0.08,5/-0.20,6/0.08,7/-0.04,8/0.12,9/-0.08,10/0.04,11/0.16,12/-0.12}
  \draw[blue!60!black,line width=1.1pt] (0.33*\k,0) -- (0.33*\k,\h);
\node[font=\scriptsize] at (2.1,1.35) {\(r_t\) 的自相关};
\end{scope}

% ---- 下右：r 平方的自相关 ----
\begin{scope}[shift={(5.7,-4.3)}]
\fill[gray!14] (0,-0.24) rectangle (4.3,0.24);
\draw[->] (0,0) -- (4.6,0) node[right,font=\scriptsize] {\(k\)};
\draw[->] (0,-0.45) -- (0,1.55);
\foreach \k/\h in {1/1.24,2/1.16,3/1.08,4/1.04,5/0.96,6/0.92,7/0.84,8/0.80,9/0.76,10/0.72,11/0.68,12/0.64}
  \draw[red!65!black,line width=1.1pt] (0.33*\k,0) -- (0.33*\k,\h);
\node[font=\scriptsize] at (2.1,1.35) {\(r_t^2\) 的自相关};
\end{scope}
\end{tikzpicture}
\end{center}

\emph{同一串数据：序列的自相关全在显著性带内，平方序列的自相关整排在带外并缓慢衰减。}

图上半部分那条固定的 \(\pm c\) 也顺带解释了告警的怪脾气。左半段几乎没有一个点触线，右半段接连穿越。阈值一动没动，穿越率却从近乎零跳到每几天一次。

\subsection{零相关不等于独立}

先把这件事算明白。设 \(r_t=\sigma_t z_t\)，其中 \(z_t\) 是均值零、方差一的独立同分布序列，\(\sigma_t>0\) 只依赖 \(t\) 时刻之前的信息。那么对任意 \(k\ge 1\)，

\[
\E[r_tr_{t-k}]
=\E\big[\E[\sigma_tz_t\given\mathcal F_{t-1}]\,r_{t-k}\big]
=\E[\sigma_tr_{t-k}\underbrace{\E[z_t]}_{=0}]
=0 .
\]

关键一步是 \(\sigma_t\) 与 \(r_{t-k}\) 都对 \(\mathcal F_{t-1}\) 可测，可以提到条件期望外面，剩下的 \(\E[z_t]=0\) 把整项抹掉。所以只要 \(z_t\) 的条件均值为零，无论 \(\sigma_t\) 的动态多复杂，一阶自相关都是零。

平方就不同了。取最简单的 \(\sigma_t^2=\omega+\alpha r_{t-1}^2\)，于是

\[
\E[r_t^2\given\mathcal F_{t-1}]=\sigma_t^2=\omega+\alpha r_{t-1}^2 ,
\]

平方序列的条件均值直接被上一期的平方牵着走。记 \(u_t=r_t^2-\sigma_t^2\)，上式可以改写成 \(r_t^2=\omega+\alpha r_{t-1}^2+u_t\)，其中 \(u_t\) 是鞅差。这就是一个关于 \(r_t^2\) 的 AR(1)，自相关是 \(\alpha^k\)。方向序列是白噪声，幅度序列是自回归——一份数据里同时住着两个模型。

零相关与独立的距离，在这里被撑到了最大：\(r_t\) 与 \(r_{t-1}\) 线性无关，但知道 \(r_{t-1}\) 会显著改变 \(r_t\) 的条件分布的宽度。相关系数只度量线性依赖，而波动率聚集是彻头彻尾的非线性依赖。

\subsection{条件方差是随机的，无条件方差可以不变}

把上面的现象写成定义。

\begin{definition}
设 \(\mathcal F_{t-1}\) 是截至 \(t-1\) 的信息（由 \(\set{r_s:s\le t-1}\) 生成的 \(\sigma\)-代数）。若
\[
\Var(r_t\given\mathcal F_{t-1})=\sigma_t^2
\]
随 \(t\) 变化且是 \(\mathcal F_{t-1}\)-可测的随机变量，则称 \(r_t\) 具有\textbf{条件异方差}。
\end{definition}

容易误读的地方是：条件异方差并不违反弱平稳。仍以 \(\sigma_t^2=\omega+\alpha r_{t-1}^2\) 为例，取无条件期望得 \(\E[\sigma_t^2]=\omega+\alpha\E[r_{t-1}^2]\)；若无条件方差存在且不随时间变化，记作 \(\bar\sigma^2\)，则

\[
\bar\sigma^2=\frac{\omega}{1-\alpha},
\qquad \alpha<1 .
\]

一个常数。均值恒为零，自协方差恒为零，弱平稳的三条要求全部满足。所以按弱平稳去筛数据，这类序列一个都筛不出来——它在二阶矩的无条件层面完全正常，异常只出现在条件层面。

后果体现在四阶矩上。若 \(z_t\) 是标准正态且 \(3\alpha^2<1\)，可以算出

\[
\frac{\E[r_t^4]}{(\E[r_t^2])^2}=\frac{3(1-\alpha^2)}{1-3\alpha^2}>3 .
\]

峰度严格大于正态的 \(3\)。换句话说，把一段条件方差变动的数据当成同方差来看，你会看到一个厚尾分布：中间更尖，两端更肥。很多``数据有厚尾''的观察，其实是``数据有条件异方差''的无条件投影。这个区分很实用——厚尾说明极端值比正态更多，条件异方差进一步说明极端值什么时候更多。

\subsection{这不只是金融的事}

收益率是最经典的例子，但它远不是唯一的。接口请求量在活动期与非活动期是两种方差；端到端延迟在容量逼近饱和时方差会先于均值爆开；模型线上错误率在数据分布漂移的窗口里剧烈抖动，漂移过去又归于安静。这些序列的共同形态是：均值稳定或缓慢变化，散布度成段落地切换。

它对日常工作的第一个后果就是开头那个告警。设阈值为 \(\pm c\)，若条件分布近似正态，则单期越界概率是

\[
P(\abs{r_t}>c\given\mathcal F_{t-1})=2\Big(1-\Phi\big(c/\sigma_t\big)\Big).
\]

这个概率对 \(\sigma_t\) 极其敏感。取 \(c=3\bar\sigma\)：当 \(\sigma_t=\bar\sigma\) 时越界概率约 \(0.27\%\)，一年不到一次；当 \(\sigma_t=2\bar\sigma\) 时变成约 \(13.4\%\)，一周好几次。波动翻一倍，误报率涨了五十倍。固定阈值假定了固定的 \(\sigma\)，而数据从来没答应过这件事。

正确的做法不是不停地调 \(c\)，而是让阈值跟着 \(\sigma_t\) 走——把告警条件改成 \(\abs{r_t}>c\,\hat\sigma_t\)，其中 \(\hat\sigma_t\) 是对条件标准差的估计。这样一来，``估计 \(\sigma_t\)'' 就从一个学术问题变成了工程需求，而这正是下一节要建的模型。

\subsection{把残差平方的自相关立成一个固定动作}

前面所有讨论都指向同一个可执行的检查。它应该出现在每一次时间序列建模的收尾环节，和\hyperref[sec:11-Mathematical-Statistics/08-Regression-and-ANOVA/01]{``经典线性回归把系数估计接到抽样分布''}里的残差检查并列，而不是等到线上出问题才回头补。

\begin{enumerate}
\tightlist
\item 先拟合条件均值模型（常数、回归、或者\hyperref[ch:052]{``ARMA 与 ARIMA：时间序列的经典模型族''}里的 ARMA），取残差 \(\hat\varepsilon_t\)。
\item 对 \(\hat\varepsilon_t\) 画自相关图，确认均值结构已经榨干——这一步大家都会做。
\item 对 \(\hat\varepsilon_t^2\) 再画一次自相关图，并对它做 Ljung--Box 检验。这一步是新增的。
\item 若要一个正式检验，做 ARCH-LM：把 \(\hat\varepsilon_t^2\) 对自己的 \(q\) 个滞后做回归，统计量 \(T R^2\) 在``没有 ARCH 效应''的原假设下渐近服从 \(\chi^2_q\)。
\end{enumerate}

第三步和第四步测的是同一件事，只是一个看图一个给 \(p\) 值。它们通过，说明残差在二阶矩上也没剩下结构，可以放心按同方差往下走；它们不通过，说明还有一层可预测性没有被建模，而这层可预测性恰好决定了区间估计、风险度量和告警阈值的可信程度。

有一点需要提前说清楚：即使检验拒绝，条件均值模型的点估计通常仍然是相合的，坏掉的是标准误。同方差假设下算出的置信区间在高波动期太窄、在平静期太宽，覆盖率不再是名义的 \(95\%\)。所以这个诊断动作的产出不是``模型错了''，而是``不确定性的刻度错了''。

下一节把 \(\sigma_t^2\) 的动态正式写下来：先是 ARCH 让条件方差依赖过去的残差平方，再是 GARCH 让它同时依赖自己的过去。
